\documentclass[11pt,a4paper]{article}
\usepackage[T1]{fontenc}
\usepackage[utf8]{inputenc}
\usepackage[czech]{babel}
\usepackage{amsmath,amssymb,amsthm}
\usepackage[margin=3cm]{geometry}
\usepackage{lmodern}
\newtheorem{theorem}{Věta}
\title{Co se ztratí, když přestaneme rozlišovat}
\author{}
\date{}
\begin{document}
\maketitle
\begin{center}
\itshape
S láskou tomu, kdo se bojí,\\
že spolu se vzpomínkou zmizí i její význam.
\end{center}
\paragraph{Model.}
Nechť $X$ je náhodná veličina s konečným nosičem
a $Y=f(X)$ její pozorování.
Shannonova entropie splňuje
\[
\boxed{H(X)=H(Y)+H(X\mid Y).}
\]
Člen $H(X\mid Y)$ vyjadřuje nejistotu o původním
stavu, která zbývá po znalosti pozorování.
\begin{proof}
Protože $Y$ je funkcí $X$, platí $H(Y\mid X)=0$.
Řetězové pravidlo tedy dává
\[
H(X)=H(X,Y)=H(Y)+H(X\mid Y).
\]
\end{proof}
\paragraph{Rozlišení a hodnota.}
Zaveďme zvlášť funkci
\[
v:\operatorname{supp}(X)\longrightarrow\mathbb R,
\]
která jednotlivým stavům přiřazuje hodnotu.
Je to další předpoklad modelu; entropie ji neurčuje.
\begin{theorem}
Z konstantního pozorování neplyne nulová hodnota:
\[
H(Y)=0
\quad\not\Longrightarrow\quad
v(X)=0.
\]
\end{theorem}
\begin{proof}
Nechť $X$ nabývá dvou hodnot $a,b$, každé
s pravděpodobností $1/2$. Položme
\[
f(a)=f(b)=y_0,
\qquad
v(a)=v(b)=1.
\]
Potom
\[
H(Y)=0,\qquad H(X\mid Y)=1,
\qquad v(X)=1.
\]
Pozorování nerozlišuje žádnou z obou možností,
ale hodnota je v obou případech nenulová.
To je protipříklad k uvedené implikaci.
\end{proof}
\paragraph{Hranice tvrzení.}
Tento příklad nedokazuje, že láska vždy přetrvá.
Dokazuje pouze, že ztráta rozlišení sama
logicky neimplikuje lhostejnost.
Lidské city nejsou v tomto modelu popsány.
\bigskip
\begin{center}
\itshape
Nevím, co zůstane po ztrátě.\\
Z toho, že něco nedokážu rozlišit,\\
však neplyne, že na ničem nezáleží.
\medskip
Matematika zde končí.\\
A člověku můžeme zůstat nablízku.
\medskip
A tak tu stojím na konci věků\\
a čekám na ostatní,\\
protože nikdo nechce být sám.
\end{center}
\end{document}
